\documentclass[11pt]{article}
% Shared preamble for per-Python-file documentation (input via \input{...}).
\usepackage[margin=1in]{geometry}
\usepackage[T1]{fontenc}
\usepackage[utf8]{inputenc}
\usepackage{lmodern}
\usepackage{hyperref}
\usepackage{xcolor}
\usepackage{listings}
\usepackage{listingsutf8}
\lstset{
  basicstyle=\ttfamily\small,
  columns=fullflexible,
  breaklines=true,
  upquote=true,
  inputencoding=utf8,
  extendedchars=true,
  frame=single,
  framerule=0.2pt,
  tabsize=2
}


\title{\texttt{model/cread.py} (Q\&A)}
\author{}
\date{}

\begin{document}
\maketitle

\paragraph{Q: What is the purpose of \texttt{model/cread.py}?}
\textbf{A:} It implements the CREAD baseline model used in this repo. The model maps feature dictionaries to a vector of per-bucket probabilities, intended to work with the discretization/reconstruction strategy implemented in \lstinline|run_cread.py|.

\paragraph{Q: What does the model output and what does it represent?}
\textbf{A:} \lstinline|forward(x_dict)| outputs a tensor of shape \lstinline|[batch, head_num]| after a sigmoid, where each head corresponds to a discretized threshold/bucket and is trained with a binary cross-entropy objective in the run script.

\paragraph{Q: How are features encoded?}
\textbf{A:} Similar to other models, it uses embeddings for sparse and sequence features (sequence pooled with a mask) and linear layers for continuous features; the pooled embeddings are fed into a shared MLP to produce a shared representation.

\paragraph{Q: How does the ``multi-head'' output work?}
\textbf{A:} After computing the shared representation, CREAD applies a separate small MLP per head (stored in \lstinline|output_mlps|), concatenates their outputs, and applies sigmoid to obtain per-head probabilities.

\paragraph{Q: Code example: multi-head construction pattern?}
\textbf{A:}
\begin{lstlisting}[language=Python]
self.output_mlps = torch.nn.ModuleList()
for _ in range(head_num):
    self.output_mlps.append(MultiLayerPerceptron(..., output_layer=True))
\end{lstlisting}

\paragraph{Q: Full source code?}
\textbf{A:}
\lstinputlisting[language=Python]{/workspace/model/cread.py}

\end{document}
